\section{Introduction}
\label{sec:introduction}

In the replay based incremental class learning, a small exemplar set represents whole training history. The backbone changes every new task, thus the exemplars selected under with old features no longer well approximate the class distributions during training. Current state of the art methods integrate selection, refresh timing, training objective, readout into one protocol, hence a performance gap between methods can not isolate the refresh rule with everything else in the method.

We propose Uniform Herding.  After each task is completed, it then divides the active budget equally among all the observed classes. And rebuilds each class's exemplar set using the greedy herding in the current feature space. Examples for reconstruction come from the active memory (bounded set of examples, retained through the tasks). Active means in the replay experience.

We then compare Uniform Herding with a faithful implementation of iCaRL and a
static replay bank on the CIFAR-100 dataset split into ten tasks. Upon arrival for iCaRL, each class was herded in a respective manner, with previous priority being maintained by truncating classes from the front that exceeded quota shrinkages per class. Each of the comparisons we make with Uniform Herding is  end-to-end, not the refresh timings but also methods for objective, readout and storage are also different. We also compare ablations of Uniform Herding on prediction rule, selection rule, distillation and head geometry, in addition to active-budget and retrieval sweeps.

At the default budgets ($M{=}2{,}000$ active exemplars, $b{=}64$ retrieval),
Uniform Herding obtains higher mean final accuracy and lower mean forgetting
than both baselines.  NME prediction and herding selection each outperform
their tested alternatives on mean accuracy; distillation primarily affects
forgetting.  The active-budget sweep changes the metrics more than the
retrieval sweep.

\paragraph{Contributions.}
\begin{itemize}
  \item We define Uniform Herding, which divides the active budget uniformly
  across classes and refreshes the exemplar set from a bounded candidate pool
  in the current representation after each task.
  \item We compare it against iCaRL and a static bank at matched active and
  retrieval budgets, with all protocol differences stated.
  \item We ablate prediction rule, selection rule, distillation, and head
  geometry, and vary active-budget and retrieval sensitivity within the
  proposed protocol.
\end{itemize}

